\section{Operational prediction and falsification map}\label{si:predictionmap}

\subsection{Independent observations and test separation}
The tests below distinguish three objects: an accounting identity, a conditional dynamical prediction, and the validity of a model--observation mapping. An identity computed from the same account inputs is not independent validation. In particular, NFA total and component entries can test arithmetic reproduction, but cannot alone identify $A_c$, $q$, conversion, or restoration. The empirical remainder $C^{\mathrm{emp}}=B-X$ is never substituted for $C^{\mathrm{mod}}=\Ups(q)A_c$ in these tests.

A physical application must first delimit a spatial system with non-overlapping land classes, record an independent land-transition matrix, and measure fast-land output per area and capital-area condition using a stated observation operator. Condition indicators such as biomass density or soil carbon need an externally justified mapping to $q$; neither indicator is assumed to be this dimensionless state. Account conversion coefficients must be fixed or estimated on calibration observations distinct from the validation interval. A failed mapping is a substantive rejection of that physical application, even when the model's internal algebra remains correct.

\begin{longtable}{>{\raggedright\arraybackslash}p{0.08\textwidth}>{\raggedright\arraybackslash}p{0.43\textwidth}>{\raggedright\arraybackslash}p{0.39\textwidth}}
\caption{Operational tests of P1--P5. Each rejection criterion is conditional on the independently checked assumptions in its row. No row is declared empirically passed by the present world-account decomposition.}\label{tab:si_predictionmap}\\
\toprule
Item & Independent data and quantities to estimate & Estimand and rejection outcome \\
\midrule
\endfirsthead
\toprule Item & Independent data and quantities to estimate & Estimand and rejection outcome \\
\midrule
\endhead
P1 & Measured converted area $v$, receiving and donor per-area annual contributions under a common valuation, and independent maintenance, retirement, restoration, condition and technology changes. Estimate $b_f$ and $\Ups(q)$ without using the held-out total $B$ to define either contribution. Prefer randomized or otherwise identified matched conversion interventions. & The conversion-only account change per area is $b_f-\Ups(q)$. A persistently opposite sign, or a held-out discrepancy exceeding propagated measurement and counterfactual uncertainty, rejects the stated mapping/mechanism after other flows are accounted for. Rising net $B$ by itself neither confirms nor rejects this marginal prediction. \\
\addlinespace
P2 & Local, matched perturbations of demand or area, repeated across operating states, with fast/capital/reserve areas and demand observed. Estimate $\mu$, $\eta_f$, $\tau_{\mathrm{conv}}$, and yield coefficients on calibration intervals; maintain the specified active-shortfall branch, fixed quality/debt and no restoration for the reduced test. & Estimate the dominant relaxation rate from the area response and compare with the predicted slow eigenvalue. Reproducible response times outside predictive uncertainty reject the calibrated reduced-rate prediction. A delay, switched regime, or changing coefficients must not be silently fitted as the same local experiment. \\
\addlinespace
P3 & A sequence of controlled or identified demand levels approaching the calibrated reduced ceiling, with repeated perturbation responses and capital-area minima. Independently check the fixed-parameter, debt-free, cap-quality, no-restoration assumptions and the one-sided active-shortfall regime. & The predicted deficit-side matrix has a demand-independent spectrum, and the equality relation reaches the area floor without a fold. An estimated relaxation rate tending to zero or a reproducible turning branch under those restrictions rejects this reduced model application. Outside those restrictions, such behavior tests a different system; finite drawdown samples cannot prove general continuity. \\
\addlinespace
P4 & For software verification, supply the exact constant history $q(s)=0.9$, $\rho_c=0.08$, the imposed cap $1$, and delays 20 or 40 yr. For a physical application, independently measure condition and identify its mapping, history, cap and regeneration rate rather than estimating them from equality of two simulated runs. & The specified projected model gives $q(t)=\min(1,0.9+0.0072t)$ and contact at $125/9$ yr. A nonvanishing integration discrepancy under refinement rejects its implementation. A mapped physical path incompatible with this law rejects the cap/history model for that system. Equality of the two constructed paths is not an empirical claim of general delay insensitivity. \\
\addlinespace
P5 & Observe capital-area losses and returns from a transition ledger; distinguish restoration from natural succession, reclassification and migration of boundaries. Compare demand reduction alone with matched return-flow interventions. Estimate conversion and return rates independently, and record the sign of $B-E$ when the policy gate acts. & With $R_{rc}=0$, the model predicts $\dot A_c=-u_c\le0$: a verified positive area return with no modeled return flow rejects the closed land-flow specification for that application. Under the declared surplus gate, positive restoration during $B<E$ rejects the gate mapping. A lower demand or later higher $B$ is not itself evidence of restoration. \\
\bottomrule
\end{longtable}

\subsection{Sampling, uncertainty and prospective rejection rules}
The validation unit, spatial boundary, calibration/validation split, intervention definition, and mapping from physical observations to model variables must be declared before inspecting validation outcomes. Measurements should resolve the fastest tested response: as a design target, spacing no larger than one-fifth of the shortest relevant conversion, demographic or regeneration timescale avoids treating an unresolved transient as a measured relaxation rate. Observation windows must also cover the delayed inputs and several predicted return times. The illustrative 1 yr conversion constant would require substantially finer observations than the annual global account series. This is a design requirement, not a claim that the present data meet it.

Parameter and measurement uncertainty should be propagated jointly, retaining covariance between area, yield and valuation. A prospective protocol can use a predeclared 95\% predictive interval for its primary held-out estimand; repeated tests require a prespecified family-wise or false-discovery adjustment. Serial observations require a dependence-aware interval or block resampling with block length chosen on calibration data. Exceedance must be replicated across the predeclared independent intervention units rather than counted repeatedly along one autocorrelated trajectory. No numerical effect-size or power guarantee follows from these illustrative parameters; pilot measurement error and an intervention effect are needed to size a field test.

For P2--P3, regime assumptions are checked independently of whether the fitted relaxation favors the prediction; cases failing those checks are reported as out-of-scope, not silently removed to improve agreement. For P1 and P5, total-account errors and land-transition errors are reported separately. Rejection of a physical mapping cannot be repaired merely by redefining the proxy remainder to equal the model capital contribution. These requirements make the proposed empirical comparison falsifiable without treating internal consistency or account arithmetic as dynamic ecological validation.
